viii Preface
with small Roman numerals, while small letters are used to label interrelated parts
in other problems. Some indication of relative difficulty is provided by a system of
starring and double starring, but there are so many criteria for judging difficulty,
and so many hints have been provided, especially for harder problems, that this
guide is not completely reliable. Many problems are so difficult, especially if the
hints are not consulted, that the best of students will probably have to attempt only
those which especially interest them; from the less difficult problems it should be
easy to select a portion which will keep a good class busy, but not frustrated. The
answer section contains solutions to about half the examples from an assortment
of problems that should provided a good test of technical competence. A separate
answer book contains the solutions of the other parts of these problems, and of all
the other problems as well. Finally, there is a Suggested Reading list, to which the
problems often refer, and a glossary of symbols.
I am grateful for the opportunity to mention the many people to whom I owe my
thanks. Jane Bjorkgren performed prodigious feats of typing that compensated for
my fitful production of the manuscript. Richard Serkey helped collect the material
which provides historical sidelights in the problems, and Richard Weiss supplied
the answers appearing in the back of the book. I am especially grateful to my
friends Michael Freeman, Jay Goldman, Anthony Phillips, and Robert Wells for
the care with which they read, and the relentlessness with which they criticized, a
preliminary version of the book. Needles to say, they are not responsible for the
deficiencies which remain, especially since I sometimes rejected suggestions which
would have made the book appear suitable for a larger group of students. I must
express my admiration for the editors and staff of W. A. Benjamin, Inc., who were
always eager to increase the appeal of the book, while recognizing the audience
for which it was intended.
The inadequacies which preliminary editions always involve were gallantly en-
dured by a rugged group of freshmen in the honors mathematics course at Brandeis
University during the academic year 1965-1966. About half of this course was
devoted to algebra and topology, while the other half covered calculus, with the
preliminary edition as the text. It is almost obligatory in such circumstances to
report that the preliminary version was a gratifying success. This is always safe—
after all, the class is unlikely to rise up in a body and protest publicly—but the
students themselves, it seems to me, deserve the right to assign credit for the thor-
oughness with which they absorbed an impressive amount of mathematics. I am
content to hope that some other students will be able to use the book to such good
purpose, and with such enthusiasm.

\textit{Waltham, Massachusetts}
\textit{February 1967}
\hfill MICHAEL SPIVAK